% Ableitung und Änderungsrate – bank/ableitung-und-aenderungsrate/ (zusammenbau v0.4, Heft)
\einheitenkopf[t1]{Ableitung und Änderungsrate}
\setcounter{aufgabe}{0}

% Anlauf (ohne Prüfkennung)
\begin{aufgabe}{Mittlere Änderungsrate – Anlauf}
\begin{teile}
\teil „Wie stark hat der Wasserstand in den ersten fünf Stunden durchschnittlich je Stunde zugenommen?“ \\ \kreuz{Zeitraum – Differenzenquotient, Sekante} \\ \kreuz{Zeitpunkt – Ableitung, Tangente}
\teil Um 8 Uhr sind 120 Besucher im Zoo, um 11 Uhr 300. Wie groß ist die mittlere Änderungsrate der Besucherzahl in diesem Zeitraum? \leerfeld Besucher je Stunde
\teil Zwei Modelle beschreiben die Abkühlung eines Getränks: $A(t) = 70 \cdot \mathrm{e}^{-0{,}1t} + 20$ und $B(t) = 90 - 3t$ ($t$ in Minuten, Temperatur in °C). Vergleiche die Beträge der mittleren Änderungsraten beider Modelle im Intervall $[0; 10]$.
\end{teile}
\end{aufgabe}

\begin{aufgabe}{Mittlere Änderungsrate – Prüfungsaufgaben}
\begin{teile}
\teil Der Graph zeigt die Anzahl $a$ der bis zum Zeitpunkt $t$ (Stunden nach 12:00 Uhr) abgegebenen Bewertungen. Lies die Werte für 13:00 Uhr und 15:00 Uhr ab und berechne, wie viele Bewertungen in diesem Zeitraum durchschnittlich je Stunde abgegeben wurden \hfill (Abitur 2025 GK).

\begin{ksys}[xmin=0,xmax=4,ymin=0,ymax=1800,xstep=1,ystep=200,xlabel=t in h,ylabel=a,ablesen]\funktion{100*\x^2}{a}\end{ksys}
\leerfeld je Stunde
\teil Der Hormonspiegel eines Patienten wird durch $h(t) = 6t \cdot \mathrm{e}^{-0{,}05t} + 40$ beschrieben ($t$ in Tagen seit Behandlungsbeginn, $h(t)$ in Prozent des Sollwerts). Berechne die mittlere Änderungsrate des Hormonspiegels in den ersten fünf Tagen \hfill (Abitur 2019 GK). \leerfeld Prozentpunkte je Tag
\teil Bei einem Laktattest beschreibt $k(x) = \frac{1}{20} \cdot (x^3 - 24x^2 + 200x - 500)$ für $10 \le x \le 16$ die Laktatkonzentration in mmol/l in Abhängigkeit von der Laufgeschwindigkeit $x$ in km/h. Berechne die mittlere Änderungsrate der Laktatkonzentration zwischen 10 und 16 km/h \hfill (Abitur 2019 GK). \leerfeld mmol/l je km/h
\teil Das Luftvolumen in der Lunge beim Einatmen wird durch $h(t) = 0{,}4t^2 - 1{,}2t + 2$ beschrieben ($t$ in Sekunden, $h(t)$ in Litern; Minimum bei $t = 1{,}5$). Gib die Bedeutung der Terme I: $h(3) - h(1{,}5)$ und II: $\frac{h(3) - h(1{,}5)}{1{,}5}$ im Sachzusammenhang an und berechne beide \hfill (Abitur 2026 GK).
\teil Ein Höhenprofil wird für $0 \le x \le 6$ durch $f(x) = (x + 2) \cdot \mathrm{e}^{-0{,}4x}$ beschrieben (1 LE = 1 km). Im Intervall $[1; 5]$ soll das Profil näherungsweise durch die Gerade durch die Punkte $(1 | f(1))$ und $(5 | f(5))$ ersetzt werden. Gib eine Gleichung dieser Geraden an \hfill (Abitur 2018 GK).
\teil Ein Höhenprofil wird durch $f(x) = (x + 2) \cdot \mathrm{e}^{-0{,}4x}$ beschrieben (1 LE = 1 km). Von einem zweiten Profil $h$ ist bekannt: $h(0) = 1{,}8$ und $h(5) = 0{,}6$. Vergleiche die Beträge der mittleren Steigungen von $f$ und $h$ im Intervall $[0; 5]$ \hfill (Abitur 2018 GK).
\end{teile}
\end{aufgabe}

\begin{aufgabe}{Mittlere Änderungsrate – Prüfungsaufgaben – weiter}
\begin{teile}
\teil Der Schienenverlauf einer Modelleisenbahn besteht aus vier gleichen Bögen von je $3{,}2$ m Länge und zwei Halbkreisen mit dem Radius $0{,}5$ m, die die Bögen verbinden. Ein Zug fährt im Durchschnitt $4$ m je Minute. Wie lange braucht er für eine volle Runde? \hfill (Abitur 2026 GK) \\ \leerfeld[min]
\teil Der Bestand eines Rohstoffs wird durch $k(t) = 0{,}1t^3 - 1{,}5t^2 + 100$ beschrieben ($t$ in Jahren, $k(t)$ in Mio. Tonnen). Die Gleichung $\frac{k(c) - k(0)}{c} = -5$ hat genau zwei Lösungen. Bestimme die kleinere Lösung und deute Gleichung und Lösung im Sachzusammenhang \hfill (Abitur 2026 GK).
\teil Die Kosten für die Produktion von $x$ Kubikmetern einer Flüssigkeit werden für $0 \le x \le 8$ durch $K(x) = x^3 - 9x^2 + 30x + 10$ beschrieben ($K(x)$ in 1\,000 Euro). Aussage: „Je größer die Produktionsmenge, desto höher sind die Kosten für einen zusätzlichen Kubikmeter.“ Beurteile die Aussage \hfill (Abitur 2018 GK).
\teil Die Kosten für die Produktion von $x$ Kubikmetern einer Flüssigkeit werden für $0 \le x \le 8$ durch $K(x) = 0{,}5x^3 - 6x^2 + 30x + 15$ beschrieben ($K(x)$ in 1\,000 Euro). Aussage: „Je größer die Produktionsmenge, desto höher sind die Kosten für einen zusätzlichen Kubikmeter.“ Beurteile die Aussage \hfill (Abitur 2018 GK).
\teil Die CO$_2$-Konzentration in einem Raum wird durch $g(x) = -500 \cdot \mathrm{e}^{-0{,}4x} + 900$ beschrieben ($x$ in Stunden, $g(x)$ in ppm). Weise nach, dass die mittlere Änderungsrate über jede Stunde $[x; x + 1]$ durch $500 \cdot (1 - \mathrm{e}^{-0{,}4}) \cdot \mathrm{e}^{-0{,}4x}$ gegeben ist, und berechne auf eine Minute genau, ab wann sie erstmals $80$ ppm je Stunde unterschreitet \hfill (Abitur 2019 GK).
\teil Die CO$_2$-Konzentration in einem Raum wird durch $g(x) = -400 \cdot \mathrm{e}^{-0{,}3x} + 800$ beschrieben ($x$ in Stunden, $g(x)$ in ppm). Weise nach, dass die mittlere Änderungsrate über jede Stunde $[x; x + 1]$ durch $400 \cdot (1 - \mathrm{e}^{-0{,}3}) \cdot \mathrm{e}^{-0{,}3x}$ gegeben ist, und berechne auf eine Minute genau, ab wann sie erstmals $60$ ppm je Stunde unterschreitet \hfill (Abitur 2019 GK).
\end{teile}
\end{aufgabe}

% Anlauf (ohne Prüfkennung)
\begin{aufgabe}{Ableitung an einer Stelle – Anlauf}
\begin{teile}
\teil $h(t)$ ist die Höhe eines Baums in cm nach $t$ Jahren. Was ist $h'(4) = 35$? \\ \kreuz{Funktionswert – Höhe oder Bestand} \\ \kreuz{Ableitungswert – Steigung oder Rate}
\teil Die Höhe eines Balls wird durch $h(t) = -5t^2 + 20t$ beschrieben ($t$ in s, $h(t)$ in m). Berechne $h'(1)$ und gib an, was die Zahl bedeutet. h'(1) = \leerfeld[m] je s
\teil $w'$ beschreibt die momentane Änderungsrate des Wasserstands einer Talsperre in cm je Stunde; $t$ ist die Zeit in Stunden nach 0:00 Uhr. Es gilt $w'(14{,}5) > 0$ und $w'(20) < 0$. Deute beide Aussagen mit Uhrzeit.
\end{teile}
\end{aufgabe}

\begin{aufgabe}{Ableitung an einer Stelle – Prüfungsaufgaben}
\begin{teile}
\teil Die Abbildung zeigt den Graphen der linearen Funktion $g$ durch $(1 | -1)$ und $(3 | 5)$. Begründe, dass $g'(3) = 3$ gilt \hfill (Abitur 2026 GK).

\begin{ksys}[xmin=-1,xmax=4,ymin=-3,ymax=7,ablesen]\gerade{3}{-4}{g}\punkt{1}{-1}{}\punkt{3}{5}{}\end{ksys}
\teil Gegeben ist $g(x) = 3\mathrm{e}^x - 3$; die Abbildung zeigt den Graphen. Berechne $g'(0)$ und veranschauliche in der Abbildung, wie man diesen Wert grafisch ermittelt \hfill (Abitur 2024 GK).

\begin{ksys}[xmin=-3,xmax=2,ymin=-4,ymax=8,ablesen]\funktion{3*exp(\x)-3}{g}\end{ksys}
\teil Gegeben ist $g(x) = 0{,}5\mathrm{e}^x + 1$; die Abbildung zeigt den Graphen. Berechne $g'(0)$ und zeichne die Tangente im Punkt $(0 | 1{,}5)$ mit einem Steigungsdreieck ein \hfill (Abitur 2024 GK).

\begin{ksys}[xmin=-3,xmax=2,ymin=-1,ymax=5,ablesen]\funktion{0.5*exp(\x)+1}{g}\end{ksys}
\teil $f$ beschreibt für $0 \le x \le 20$ die momentane Änderungsrate des Inhalts eines Tanks in kg je Stunde ($x$ Zeit in Stunden seit Beobachtungsbeginn); der Punkt $(3 | 180)$ liegt auf dem Graphen von $f$. Gib die Bedeutung der Koordinaten im Sachzusammenhang an \hfill (Abitur 2021 GK).
\teil Die Schaufel eines Wasserrads wird durch die Graphen von $f(x) = 0{,}1x^3 - x^2 + 2x$ und $p(x) = -x^2 + 3x - 0{,}5$ begrenzt. Weise nach, dass $p'(0{,}5) = f'(0)$ gilt, und erläutere die anschauliche Bedeutung im Sachzusammenhang \hfill (Abitur 2021 GK).
\end{teile}
\end{aufgabe}

\begin{aufgabe}{Ableitung an einer Stelle – Prüfungsaufgaben – weiter}
\begin{teile}
\teil Ein Logo wird durch die Graphen von $u(x) = 0{,}2x^3$ (untere Linie) und $v(x) = x^2 - 0{,}1x^3$ (obere Linie) begrenzt. Aussage: „Für jedes $x$ aus $]0; 3[$ ist die Steigung des Graphen von $v$ größer als die des Graphen von $u$.“ Beurteile die Aussage \hfill (Abitur 2020 GK).
\teil Zwei Rampen werden durch $u(x) = 0{,}5x^2$ und $v(x) = 3x - 0{,}25x^2$ beschrieben. Aussage: „Für jedes $x$ aus $]0; 4[$ ist die Rampe $v$ steiler als die Rampe $u$.“ Beurteile die Aussage \hfill (Abitur 2020 GK).
\end{teile}
\end{aufgabe}

% Anlauf (ohne Prüfkennung)
\begin{aufgabe}{Von der Sekante zur Tangente – Anlauf}
\begin{teile}
\teil Für eine Funktion $f$ ist der Term $\frac{f(5) - f(2)}{5 - 2}$ gegeben. Was gibt er an? \\ \kreuz{Sekantensteigung – mittlere Änderungsrate} \\ \kreuz{Tangentensteigung – momentane Änderungsrate}
\teil Gegeben ist $f(x) = x^2$ und der Punkt $P(2 | 4)$ auf dem Graphen. Berechne die Steigung $m$ der Sekante durch $P$ und $Q(2 + h | f(2 + h))$ für $h = 1$; $0{,}5$; $0{,}1$; $0{,}01$ und trage sie in die Tabelle ein. Welchem Wert nähern sich die Steigungen? Das ist die Tangentensteigung $f'(2)$.

\wertetabelle{h}{m}{1,0{,}5,0{,}1,0{,}01}
f'(2) = \leerfeld
\teil Die Abbildung zeigt den Graphen von $g(x) = -0{,}5x^3 + 1{,}5x + 3$ mit dem Wendepunkt $W(0 | 3)$. Berechne die mittlere Steigung des Graphen im Bereich $-2 \le x \le 2$. Bestimme grafisch die Steigung des Graphen im Wendepunkt.

\begin{ksys}[xmin=-3,xmax=3,ymin=-1,ymax=7,ablesen]\funktionab{-0.5*\x^3+1.5*\x+3}{g}{-2.4}{2.4}\end{ksys}
\end{teile}
\end{aufgabe}

\begin{aufgabe}{Von der Sekante zur Tangente – Prüfungsaufgaben}
\begin{teile}
\teil Die Temperatur eines Kaffees wird durch $T(x) = 55 \cdot \mathrm{e}^{-x/200} + 21$ beschrieben ($x$ in Minuten seit dem Einschenken, $T(x)$ in °C); $T'(x) = -0{,}275 \cdot \mathrm{e}^{-x/200}$. Prüfe, ob die mittlere Änderungsrate der Temperatur in den ersten 40 Minuten um mehr als $10\,\%$ von der momentanen Änderungsrate nach 40 Minuten abweicht \hfill (Abitur 2026 GK).
\teil Die Anzahl der Bakterien in einer Kultur wird durch $B(t) = 200 \cdot \mathrm{e}^{0{,}3t}$ beschrieben ($t$ in Stunden, $B(t)$ in Tausend); $B'(t) = 60 \cdot \mathrm{e}^{0{,}3t}$. Prüfe, ob die mittlere Änderungsrate der ersten zwei Stunden um mehr als $25\,\%$ von der momentanen Änderungsrate nach zwei Stunden abweicht \hfill (Abitur 2026 GK).
\teil Gegeben ist $f(x) = x^3 - 3x$. Weise nach, dass mindestens eine Stelle $x_0$ die Gleichung $f'(x_0) = \frac{f(3) - f(0)}{3 - 0}$ erfüllt, und gib die geometrische Bedeutung von $x_0$ an \hfill (Abitur 2022 GK).
\teil Gegeben sind $f(x) = \sqrt{x}$ für $x \ge 0$ und die Gerade $g$: $y = 0{,}5x$. Der Graph von $f$ und $g$ schneiden sich in zwei Punkten. Im Intervall zwischen den x-Koordinaten der Schnittpunkte gibt es eine Stelle, an der die lokale Änderungsrate von $f$ mit der mittleren Änderungsrate von $f$ in diesem Intervall übereinstimmt. Bestimme diese Stelle \hfill (Abitur 2024 GK). x = \leerfeld
\end{teile}
\end{aufgabe}

% Anlauf (ohne Prüfkennung)
\begin{aufgabe}{Die Rate als Funktion – Anlauf}
\begin{teile}
\teil $f(t)$ ist die Wassermenge in einem Becken in Litern nach $t$ Minuten. Was beschreibt $f$, und wo ist die größte Zuflussrate zu suchen? \\ \kreuz{Bestand (Höhe, Anzahl, Menge) – größte Rate bei den Nullstellen der zweiten Ableitung} \\ \kreuz{Rate (je Zeiteinheit) – größte Rate bei den Nullstellen der ersten Ableitung}
\teil Die Zuflussrate in ein Becken wird durch $r$ beschrieben ($t$ in Stunden, $r(t)$ in Litern je Stunde); es gilt $r'(t) = (6 - t) \cdot \mathrm{e}^{-0{,}5t}$. Gib die Nullstelle von $r'$ an und nenne den Zeitpunkt der größten Zuflussrate an. t = \leerfeld
\teil Die Temperatur in einem Gewächshaus wird durch $T(t) = -0{,}1t^3 + 1{,}8t^2 + 5$ beschrieben ($t$ in Stunden nach Sonnenaufgang, $T(t)$ in °C, $0 \le t \le 12$). Wann steigt die Temperatur am schnellsten, und wie groß ist die Änderungsrate dann? t = \leerfeld , \leerfeld °C je h
\end{teile}
\end{aufgabe}

\begin{aufgabe}{Die Rate als Funktion – Prüfungsaufgaben}
\begin{teile}
\teil Gegeben ist $f(x) = 0{,}5x^3 - 3x + 2$. Unter den Tangenten an den Graphen von $f$ hat eine die kleinste Steigung. Berechne diese Steigung \hfill (Abitur 2019 GK). m = \leerfeld
\teil Ein Höhenprofil wird für $0 \le x \le 8$ durch $p(x) = (2x + 1) \cdot \mathrm{e}^{-0{,}4x}$ beschrieben (1 LE = 1 km); $p'(x) = (1{,}6 - 0{,}8x) \cdot \mathrm{e}^{-0{,}4x}$. Weise nach, dass es im Intervall $[2; 8]$ eine Stelle gibt, an der die Steigung des Profils kleiner als $-0{,}33$ ist \hfill (Abitur 2018 GK).
\teil Die Höhe eines Bambus wird durch $w(t) = -t^3 + 12t^2$ beschrieben ($t$ in Wochen, $w(t)$ in cm, $0 \le t \le 8$); $w'(t) = -3t^2 + 24t$. Berechne die höchste Wachstumsgeschwindigkeit in cm je Woche; die notwendige Bedingung genügt \hfill (Abitur 2019 GK). \leerfeld[cm] je Woche
\teil Die momentane Förderrate einer Mine wird durch $k'(x) = 0{,}3x \cdot \mathrm{e}^{-0{,}1x}$ beschrieben ($x$ in Jahren seit Beginn der Förderung, $k'(x)$ in Mio. t je Jahr). Berechne den Zeitpunkt der größten Förderrate und ihren Wert \hfill (Abitur 2026 GK). x = \leerfeld , \leerfeld Mio. t je Jahr
\teil Die Höhe einer Drohne wird für $0 \le t \le 15$ durch $h(t) = 0{,}2t^3 - 6t^2 + 45t$ beschrieben ($t$ in Minuten, $h(t)$ in m); die Phase des Aufstiegs dauert von 0 bis 5 Minuten, $h'(t) = 0{,}6t^2 - 12t + 45$. Bestimme den Zeitraum in dieser Phase, in dem die Steiggeschwindigkeit mindestens $23{,}4$ m je min beträgt \hfill (Abitur 2022 GK).
\teil Der Füllstand eines Tanks wird für $0 \le t \le 10$ durch $f(t) = -0{,}25t^3 + 3t^2 + 5$ beschrieben ($t$ in Minuten, $f(t)$ in Litern); die Zuflussphase dauert von 0 bis 8 Minuten, $f'(t) = -0{,}75t^2 + 6t$. Bestimme den Zeitraum in dieser Phase, in dem die Zuflussrate mindestens $9$ L je min beträgt \hfill (Abitur 2022 GK).
\end{teile}
\end{aufgabe}

\begin{aufgabe}{Die Rate als Funktion – Prüfungsaufgaben – weiter}
\begin{teile}
\teil Zwei Bäume wachsen nach $a(t) = 18t^2 - t^3$ und $b(t) = 3t^2$ ($t$ in Jahren, Höhe in cm, $0 \le t \le 12$). Berechne den Zeitpunkt, zu dem der Unterschied der Wachstumsgeschwindigkeiten am größten ist, und diesen Unterschied \hfill (Abitur 2019 GK).
\teil Zwei Fahrzeuge starten gleichzeitig; ihre Wege sind $s(t) = 0{,}05t^3 + 2t$ und $u(t) = 1{,}5t^2$ ($t$ in Sekunden, Weg in m, $0 \le t \le 15$). Berechne den Zeitpunkt, zu dem der Unterschied ihrer Geschwindigkeiten am größten ist, und diesen Unterschied \hfill (Abitur 2019 GK).
\teil Die Verkaufszahlen zweier Produkte werden durch $a(t) = t^3 - 9t^2 + 30t$ und $b(t) = 15t$ beschrieben ($t$ in Monaten, verkaufte Stück in Tausend, $0 \le t \le 6$). Bestimme die Zeitpunkte, zu denen der Unterschied der Verkaufsraten am größten ist \hfill (Abitur 2019 GK).
\teil Die Temperatur eines Getränks im Kühlschrank wird durch $T(x) = 50 \cdot \mathrm{e}^{-0{,}02x} + 6$ beschrieben ($x$ in Minuten, $T(x)$ in °C); $T(x) - 6$ ist die Differenz zur Kühlschranktemperatur, $T'(x) = -\mathrm{e}^{-0{,}02x}$. Aussage: Es gibt eine Zahl $c$, sodass für alle $x$ gilt $T(x) - 6 = c \cdot T'(x)$. Begründe, dass die Aussage im Modell wahr ist \hfill (Abitur 2026 GK).
\teil Die Menge eines Medikaments im Blut wird durch $m(t) = 80 \cdot \mathrm{e}^{-0{,}25t}$ beschrieben ($t$ in Stunden, $m(t)$ in mg); $m'(t) = -20 \cdot \mathrm{e}^{-0{,}25t}$. Aussage: Zu jedem Zeitpunkt ist die Menge proportional zur momentanen Abbaurate, $m(t) = c \cdot m'(t)$ für alle $t$. Begründe, dass die Aussage im Modell wahr ist \hfill (Abitur 2026 GK).
\teil Der Ladestand eines Akkus wird durch $L(t) = 100 - 60 \cdot \mathrm{e}^{-0{,}05t}$ beschrieben ($t$ in Minuten, $L(t)$ in Prozent); $L'(t) = 3 \cdot \mathrm{e}^{-0{,}05t}$. Aussage: Es gibt eine Zahl $c$, sodass für alle $t$ gilt $L(t) - 100 = c \cdot L'(t)$. Begründe, dass die Aussage im Modell wahr ist \hfill (Abitur 2026 GK).
\teil Die Einwohnerzahl einer Stadt wird durch $B(t) = 500 \cdot \mathrm{e}^{0{,}04t} + 200$ beschrieben ($t$ in Jahren, $B(t)$ in Tausend); $B'(t) = 20 \cdot \mathrm{e}^{0{,}04t}$. Aussage: Es gibt eine Zahl $c$, sodass für alle $t$ gilt $B(t) - 200 = c \cdot B'(t)$. Begründe, dass die Aussage im Modell wahr ist, und gib $c$ an \hfill (Abitur 2026 GK). c = \leerfeld
\end{teile}
\end{aufgabe}
